\subsection{Тракт данных и процедуры калибровки}
\label{sec:calibration}

\subsubsection{Сквозной тракт данных}

Полный цикл получения радиолокационного изображения включает
последовательность преобразований, представленную на
рис.~\ref{fig:data-pipeline}.

\begin{figure}[H]
    \centering
    \fbox{\parbox[c][0.22\textheight][c]{0.88\textwidth}{\centering
    Место для схемы сквозного тракта данных:\\[4pt]
    Задание $T$, $I$ $\to$ Оптическое гетеродинирование $\to$
    $f_{\mathrm{RF}}$ $\to$\\
    Облучение объекта $\to$ Приём отражённого сигнала $\to$
    Синхронное детектирование ($I_k$, $Q_k$) $\to$\\
    Оцифровка АЦП $\to$ Формирование
    $S_{21}(f_k)$ $\to$ ОБПФ $\to$ A-скан $\to$ B-скан.}}
    \caption{Сквозная схема тракта данных от генерации сигнала до
    радиолокационного изображения}
    \label{fig:data-pipeline}
\end{figure}

Цикл начинается с \textbf{управляющего сигнала}: ПО задаёт температуру и/или
ток лазера, определяя текущую частоту $f_k$ радиофотонного генератора. На
выходе фотодиода формируется монохроматический \textbf{СВЧ-сигнал} на частоте
$f_k$, который через передающую антенну облучает объект. Принятый
\textbf{аналоговый отклик} поступает на синхронный детектор, выдающий пару
напряжений $U_I$ и $U_Q$, пропорциональных $I_k$ и $Q_k$. АЦП преобразует
их в 16-битные целые числа, при стробированном захвате выполняя усреднение
по 1000 окнам согласно~\eqref{eq:window-avg}. После нормировки на амплитуду
передаваемого сигнала формируется \textbf{комплексный S-параметр}
$S_{21}(f_k) \propto I_k + iQ_k$, а совокупность значений
$\{S_{21}(f_k)\}_{k=0}^{N-1}$ составляет один полный \textbf{частотный свип}.
Обратное преобразование Фурье этого свипа даёт \textbf{A-скан} --- отклик во
временной области. Последовательность A-сканов, полученных при перемещении
антенного блока вдоль профиля, формирует \textbf{B-скан} --- двумерное
радиолокационное изображение.

\subsubsection{Временна\'{я} диаграмма одного частотного свипа}

На каждом частотном шаге $k$ (рис.~\ref{fig:timing-diagram}) управляющее ПО
передаёт новое значение тока или температуры на плату управления лазерами.
После паузы стабилизации (20--100\,мкс при токовой или 3--10\,мс при
температурной перестройке) стробирующий сигнал 20\,кГц открывает окно захвата
длительностью $1/20\,000 = 50$\,мкс, в пределах которого АЦП выполняет
100 отсчётов с тактовой частотой 2\,МГц. Процедура повторяется $M = 1000$ раз,
результаты усредняются, после чего усреднённые значения $I_k$, $Q_k$
передаются на ПК и сохраняются, а система переходит к следующему частотному
шагу $f_{k+1}$.

\begin{figure}[H]
    \centering
    \fbox{\parbox[c][0.18\textheight][c]{0.85\textwidth}{\centering
    Место для временно\'{й} диаграммы одного свипа:\\[4pt]
    $f_k \to$ стабилизация $\to$ 1000 окон $\times$ 100 отсчётов
    $\to$ усреднение $\to$ передача $\to$ $f_{k+1}$.}}
    \caption{Временна\'{я} диаграмма измерительного цикла для одного
    частотного шага}
    \label{fig:timing-diagram}
\end{figure}

Полное время одного свипа зависит от числа частотных шагов $N$ и задержки
стабилизации на каждом шаге. При токовой перестройке с минимальной задержкой
$\tau = 20$\,мкс и $N = 1000$ шагов время свипа составляет порядка единиц
секунд (с учётом накопления 1000 окон на каждом шаге). При температурной
перестройке время увеличивается пропорционально задержке стабилизации.

\subsubsection{Процедуры калибровки}

Для получения количественно корректных результатов перед началом
измерительной серии выполняется набор калибровочных процедур.

\paragraph{Калибровка частотной оси.}
Зависимость между управляющим параметром (номер шага перестройки $x$) и
фактической частотой СВЧ-сигнала в общем случае нелинейна. Для её описания
используется полиномиальная аппроксимация:
\begin{equation}
    f_{\mathrm{cal}}(x) = c_0 + c_1 x + c_2 x^2,
    \label{eq:freq-cal}
\end{equation}
где коэффициенты $c_0$, $c_1$, $c_2$ определяются экспериментально путём
измерения частоты биения анализатором спектра при нескольких значениях
управляющего параметра. После определения коэффициентов все шаговые индексы
пересчитываются в физические частоты, и данные при необходимости
интерполируются на равномерную частотную сетку.

\paragraph{Калибровка в свободном пространстве (open-air).}
Перед каждой серией измерений выполняется регистрация полного частотного
свипа при направлении антенн в свободное пространство (без отражающего
объекта). Полученный отклик $S_{\mathrm{oa}}(f_k)$ содержит систематические
артефакты: прямую связь между антеннами, отражения от элементов конструкции,
собственные резонансы тракта. При обработке данных этот референсный свип
вычитается из измеренного:
\begin{equation}
    \Delta S(f_k) = S(f_k) - S_{\mathrm{oa}}(f_k),
    \label{eq:open-air-subtraction}
\end{equation}
что соответствует процедуре формирования приведённого S-параметра, описанной
в теоретической части.

\paragraph{Калибровка нуля времени.}
После обратного преобразования Фурье начало временно\'{й} оси ($t = 0$) может
не совпадать с плоскостью антенн из-за электрических задержек в тракте. Для
определения поправки $t_0$ выполняется измерение с металлическим отражателем,
расположенным на известном расстоянии $d_0$ от антенн:
\begin{equation}
    t_0 = t_{\mathrm{meas}} - \frac{2d_0}{c},
    \label{eq:t0-calibration}
\end{equation}
где $t_{\mathrm{meas}}$ --- положение максимума отражения на A-скане.
Полученная поправка применяется ко всем последующим измерениям.

\paragraph{Амплитудная калибровка.}
Неравномерность амплитудно-частотной характеристики СВЧ-тракта
компенсируется путём нормировки на калибровочный свип. При первом способе
(\textbf{проекция между огибающими}) из калибровочного свипа извлекаются
верхняя и нижняя огибающие с помощью скользящего среднего по локальным
экстремумам, и измеренный свип линейно проецируется в стандартизованный
диапазон. При втором способе (\textbf{поэлементное деление}) каждый отсчёт
просто делится на соответствующее значение калибровочного свипа.

\paragraph{Комплексная (I/Q) калибровка.}
При фазочувствительных измерениях калибровка выполняется в комплексной
форме. Регистрируется калибровочный свип $C(f_k) = I_{\mathrm{cal}}(f_k) +
iQ_{\mathrm{cal}}(f_k)$ при известных условиях, после чего последующие
измерения нормируются поэлементным комплексным делением:
\begin{equation}
    S_{21}^{\mathrm{norm}}(f_k) = \frac{S_{21}(f_k)}{C(f_k)},
    \label{eq:complex-cal}
\end{equation}
что одновременно компенсирует как амплитудную, так и фазовую неравномерность
тракта.

\subsubsection{Средства верификации и тестирования}

Для обеспечения надёжности и воспроизводимости измерений в состав
программного комплекса входят \textbf{утилита проверки тактовой частоты},
подсчитывающая число тактов АЦП между фронтами стробирующего сигнала для
контроля стабильности внешнего тактирования; \textbf{эмулятор виртуального
порта}, воспроизводящий ранее записанные данные с регулируемой скоростью и
позволяющий отлаживать алгоритмы без подключения аппаратуры; а также наборы
\textbf{автоматических тестов} --- 75 тестов для модуля управления лазерами
(валидация параметров, кодирование/декодирование протокола, сквозное
взаимодействие) и собственный набор тестов модуля визуализации (вычисление
БПФ, нормировка, калибровка, детекция пиков).

Совокупность описанных калибровочных процедур и средств верификации
формирует воспроизводимый измерительный протокол, позволяющий перейти от
аппаратного уровня радиофотонного радара к получению количественно
интерпретируемых радиолокационных данных.
